\section{核能基础：裂变、聚变与可控}

本章先消化基础概念。核能的共同来源是质量亏损：反应前后的总质量不同，少掉的质量通过爱因斯坦质能方程转化为能量。裂变是重原子核分裂成较轻原子核，聚变是轻原子核合并成较重原子核；两者都可以释放能量，但燃料、反应条件、废料和安全特征完全不同。

核聚变为什么必须加上“可控”两个字？因为不可控释放就是氢弹一类的爆炸过程；可控核聚变要做的是按人类需要的功率、时间和安全边界稳定输出能量。这就是从“释放能量”到“能源产品”的巨大差别：前者证明反应能发生，后者要求反应可维持、可调节、可取热、可并网、可维护。

\[
E=\Delta m c^2
\]

其中，$E$ 表示释放的能量，$\Delta m$ 表示反应前后减少的质量，$c$ 表示光速。因为 $c^2$ 极大，所以很小的质量差也能对应巨大能量。这个公式解释了核能的能量密度，也解释了为什么核聚变一旦可控，会被称为能源圣杯。

\begin{knowledgebox}{术语消化：裂变、聚变、可控聚变}
\begin{tabularx}{\textwidth}{p{0.18\textwidth}p{0.34\textwidth}X}
\toprule
概念 & 核心机制 & 工程含义 \\
\midrule
裂变 & 重核分裂成较轻核，释放质量差对应的能量 & 已经商业化，但有放射性燃料、废料和安全监管压力。 \\
聚变 & 轻核碰撞并融合成较重核，释放能量 & 燃料潜力巨大，但需要极高温度和约束条件。 \\
不可控聚变 & 能量在极短时间内释放 & 武器逻辑，不是发电逻辑。 \\
可控核聚变 & 按设计功率稳定释放能量 & 需要约束、控制、取热、材料、运维和电网系统共同过线。 \\
\bottomrule
\end{tabularx}
\end{knowledgebox}

\begin{figure}[H]
\centering
\includegraphics[width=0.96\textwidth]{fusion-energy-chain.png}
\caption{可控核聚变能量链：从燃料、等离子体到电网输出的核心路径。自制概念图，依据 00:04:22--00:08:34 对谈内容整理。}
\end{figure}

\begin{knowledgebox}{读图：核聚变电站不是“点燃”等离子体就结束}
图中 Fuel 到 Plasma 是物理反应入口；Confinement 是让高温等离子体不碰到固体壁；Heat 和 Grid 则代表工程产品化：能量必须被取出来、转成电、稳定接入电网。很多讨论只停在前两步，但商业化电站必须走完整条链。
\end{knowledgebox}

\subsection{本章小结}

裂变和聚变都来自质量亏损，但可控核聚变的挑战在于“稳定、可控、低成本地释放并利用能量”。它不是一次爆炸，而是一套长期运行的能源系统。

